\providecommand*{\olmrsourcecitation}[1]{%
  \ifcsname olmrsourcecitation@#1\endcsname
    \csname olmrsourcecitation@#1\endcsname
  \else
    \PackageError{slc-mr}{Unknown scholarly citation #1}{}%
  \fi
}
\expandafter\newcommand\csname olmrsourcecitation@cantor-treatment\endcsname{\footnote{ऑलिव्हर डाइझर, \emph{Einführung in die Mengenlehre}, पृ.~218, आणि लिओपोल्डिनाचे \href{https://www.leopoldina.org/en/members/member-list/detail/georg-cantor}{कँटर यांचे सदस्यचरित्र} पाहा.}}
\expandafter\newcommand\csname olmrsourcecitation@gentzen-publication\endcsname{\footnote{गेंटझेन यांच्या लेखांच्या \href{https://doi.org/10.1007/BF01201353}{पहिल्या} आणि \href{https://doi.org/10.1007/BF01201363}{दुसऱ्या} भागांच्या प्रकाशकनोंदी पाहा.}}
\expandafter\newcommand\csname olmrsourcecitation@godel-degree\endcsname{\footnote{व्हिएन्ना विद्यापीठाचे \href{https://geschichte.univie.ac.at/de/artikel/von-der-berechenbarkeit-zum-modernen-computer}{इतिहासवर्णन} आणि इन्स्टिट्यूट फॉर ॲडव्हान्स्ड स्टडीची \href{https://www.ias.edu/scholars/godel}{गोडेल यांची नोंद} पाहा.}}
\expandafter\newcommand\csname olmrsourcecitation@russell-prison-and-pacifism\endcsname{\footnote{तुरुंगवास आणि सुटकेच्या दिनांकासाठी मॅकमास्टर विद्यापीठाचा \href{https://russell-letters.mcmaster.ca/chronology}{कालक्रम} पाहा. बलप्रयोगावरील रसेल यांची भूमिका त्यांच्या \href{https://bracers.mcmaster.ca/30909}{पत्रात} स्पष्ट होते.}}
\expandafter\newcommand\csname olmrsourcecitation@turing-bombe\endcsname{\footnote{बॉम्ब यंत्राविषयी राष्ट्रीय संगणन संग्रहालयाची \href{https://www.tnmoc.org/bombe}{नोंद} पाहा.}}
\expandafter\newcommand\csname olmrsourcecitation@turing-clarke\endcsname{\footnote{क्लार्क यांच्यावरील जीसीएचक्यूचे \href{https://www.gchq.gov.uk/information/joan-clarke}{चरित्र} पाहा.}}
\expandafter\newcommand\csname olmrsourcecitation@zermelo-appointments\endcsname{\footnote{झर्मेलो यांच्या प्राध्यापकपदाचा कालक्रम आणि राजीनाम्याची पार्श्वभूमी फ्रायबुर्ग विद्यापीठाच्या \href{https://uni-freiburg.de/math/geschichte-i/}{विभागाच्या इतिहासात} आणि \href{https://home.mathematik.uni-freiburg.de/mildenberger/freiburg2014/zermelo.html}{झर्मेलो यांच्या चरित्रात} दिली आहे.}}
\expandafter\newcommand\csname olmrsourcecitation@avigad-logic-and-computation\endcsname{या अध्यायाची ओपन लॉजिक प्रकल्पातील मांडणी जेरेमी अविगॅड यांच्या \href{https://www.andrew.cmu.edu/user/avigad/Teaching/landc_notes.pdf\#page=105}{\emph{Logic and Computation}} या व्याख्यानांच्या टिपणांतील प्रकरण ८ वर आधारित आहे (जानेवारी २००२, पृ.~९९ पुढे).}
